% !TEX program = pdflatex
% =====================================================================
%  Informe técnico PID IV — Alain
%  Réplica en LaTeX del documento Word original.
%  Compilar con LuaLaTeX o XeLaTeX (usa Arial y Times New Roman del sistema)
%  o con pdfLaTeX (usa Helvetica, de métricas idénticas a Arial):
%      latexmk -lualatex main.tex      |  latexmk -pdf main.tex
%
%  Medidas tomadas del .docx original:
%   - Papel Carta (21,59 x 27,94 cm), márgenes laterales 2 cm,
%     encabezado y pie a 2 cm del borde; márgenes superior/inferior
%     propios de cada sección.
%   - Arial; interlineado 1,5 (Word: 1,5 x 1,15 x tamaño).
%   - Sin división silábica automática, como en Word.
% =====================================================================
\documentclass[12pt]{article}

\usepackage[paperwidth=612bp,paperheight=792bp,
            left=56.7bp,right=56.7bp,top=82.35bp,bottom=87.19bp,
            headheight=13bp,headsep=14.15bp,footskip=28.37bp]{geometry}
\usepackage{iftex}
\ifPDFTeX
  % pdfLaTeX: Helvetica (Nimbus Sans) tiene exactamente las mismas
  % anchuras que Arial, por lo que los saltos de línea coinciden.
  \usepackage[T1]{fontenc}
  \usepackage[utf8]{inputenc}
  \usepackage{lmodern}
  \usepackage{helvet}
  \renewcommand{\familydefault}{\sfdefault}
  \newcommand\timesnr{\fontfamily{ptm}\selectfont}
\else
  % LuaLaTeX / XeLaTeX: fuentes originales del documento
  \usepackage{fontspec}
  \setmainfont{Arial}
  \setsansfont{Arial}
  \newfontfamily\timesnr{Times New Roman}
\fi
\usepackage[spanish,es-noshorthands,es-lcroman,es-nodecimaldot]{babel}
\usepackage{graphicx}
\usepackage{xcolor}
\usepackage{colortbl}
\usepackage{longtable}
\usepackage{csvsimple-l3}
\usepackage[normalem]{ulem}
\usepackage{fancyhdr}
\usepackage{enumitem}
\usepackage[hidelinks,hypertexnames=false,bookmarksnumbered=false,
            pdftitle={Sistema de microtransacciones mediante plataformas de pago nacionales},
            pdfauthor={Alain},
            pdfsubject={Informe técnico. Proyecto Investigación y Desarrollo IV}]{hyperref}

% ---------------------------------------------------------------------
%  Párrafo base (estilo "Normal" de Word)
% ---------------------------------------------------------------------
\setlength{\parindent}{0pt}
\setlength{\parskip}{0pt}
\raggedbottom                     % Word no estira el espacio vertical
\widowpenalty=10000 \clubpenalty=10000 % control de viudas y huérfanas de Word
\hyphenpenalty=10000              % Word no divide palabras
\tolerance=2000 \emergencystretch=3em \hbadness=10000 \vbadness=10000
% Justificado al estilo de Word: los espacios se estiran libremente pero
% solo se comprimen hasta un 25 % (TeX permite por defecto mucho más).
\AddToHook{selectfont}{\spaceskip=\fontdimen2\font
  plus \fontdimen3\font minus .25\fontdimen2\font\relax}

% Tamaños con interlineado de Word (factor de línea de Arial = 1,15)
\newcommand\doceUnoCinco{\fontsize{12bp}{20.7bp}\selectfont}      % 12 pt, 1,5
\newcommand\onceUnoCinco{\fontsize{11bp}{18.975bp}\selectfont}    % 11 pt, 1,5
\newcommand\onceUnoQuince{\fontsize{11bp}{14.5475bp}\selectfont}  % 11 pt, 1,15
\newcommand\onceSencillo{\fontsize{11bp}{12.65bp}\selectfont}     % 11 pt, sencillo
\newcommand\doceSencillo{\fontsize{12bp}{13.8bp}\selectfont}      % 12 pt, sencillo
\newcommand\diezSencillo{\fontsize{10bp}{11.5bp}\selectfont}      % 10 pt, sencillo

\renewcommand\normalsize{\fontsize{12bp}{20.7bp}\selectfont}

% Párrafo vacío de Word (una línea en blanco del tamaño actual)
\newcommand\vacio{\leavevmode\par}

% Espacio "anterior" de un párrafo situado al inicio de página
\newcommand\espacioInicial[1]{\vspace*{\dimexpr#1-\topskip\relax}\prevdepth=-1000pt\relax}

% Versalitas sintéticas, como las genera Word (minúsculas al 80 %)
\ExplSyntaxOn
\NewDocumentCommand\versalitas{m}
  {
    \text_map_inline:nn {#1}
      {
        \str_if_eq:eeTF { \text_uppercase:n {##1} } { ##1 }
          { ##1 }
          { { \fontsize{\fp_eval:n {0.8 * \use:c{f@size}} pt}{0pt}\selectfont \text_uppercase:n {##1} } }
      }
  }
\ExplSyntaxOff

% ---------------------------------------------------------------------
%  Secciones de Word: márgenes, encabezado y pie
% ---------------------------------------------------------------------
%  \seccion{superior}{sep. encabezado}{inferior}{sep. pie}{primera línea}
%  El encabezado empieza a 2 cm del borde superior y el pie termina a 2 cm
%  del borde inferior. Como en Word, la línea base de la primera línea queda
%  a la altura de la ascendente de Arial (0,9375 x tamaño) bajo el margen y
%  el espacio adicional del interlineado se añade debajo del texto; el margen
%  inferior incluye la descendente (2,54 pt) para que una línea solo quepa
%  si su descendente no rebasa el margen de Word.
\newcommand\seccion[5]{%
  \newgeometry{left=56.7bp,right=56.7bp,top=#1,bottom=#3,
               headheight=13bp,headsep=#2,footskip=#4}%
  \setlength{\topskip}{#5}}

\newcommand\textoEncabezado{}
\fancypagestyle{informe}{%
  \fancyhf{}%
  \fancyhead[R]{\diezSencillo\itshape\rule[-2.12bp]{0pt}{11.5bp}\textoEncabezado}%
  \fancyfoot[C]{\diezSencillo\thepage}%
  \renewcommand{\headrulewidth}{0pt}%
  \renewcommand{\footrulewidth}{0pt}}
\fancypagestyle{portada}{\fancyhf{}%
  \renewcommand{\headrulewidth}{0pt}\renewcommand{\footrulewidth}{0pt}}

% ---------------------------------------------------------------------
%  Títulos (estilos "Título 1" y "Título 2" de Word)
% ---------------------------------------------------------------------
%  Título 1: Arial 14 negrita, versalitas, justificado, 1,5,
%  11,35 pt antes y después; siempre al inicio de página.
%  \tituloI*{...} : sin versalitas (así está "Introducción" en el original).
%  \tituloI[sangría]{...} : "Anexos" lleva una numeración vacía seguida de
%                            tabulación, que lo desplaza 1,27 cm.
\NewDocumentCommand\tituloI{s O{0pt} m}{%
  \clearpage
  \espacioInicial{11.35bp}%
  \phantomsection
  \addcontentsline{toc}{section}{\texorpdfstring{\MakeUppercase{#3}}{#3}}%
  {\parskip=0pt\fontsize{14bp}{24.15bp}\selectfont\bfseries
   \hspace*{#2}\rule[-2.97bp]{0pt}{16.095bp}\IfBooleanTF{#1}{#3}{\versalitas{#3}}\par}%
  \nopagebreak
  % 11,35 pt después + resto de la línea de 14 pt (24,15 - 13,125) hasta la
  % ascendente de la primera línea del texto que sigue
  \vspace{\dimexpr 11.35bp + 11.025bp - \baselineskip
          + 0.9375\dimexpr\csname f@size\endcsname pt\relax - \parskip\relax}}

%  Título 2: Arial 12 negrita, justificado, 1,5, 5,65 pt antes y después.
\newcommand\tituloII[1]{%
  \par\vspace{5.65bp}%
  {\parskip=0pt\doceUnoCinco\bfseries
   \leavevmode\phantomsection\addcontentsline{toc}{subsection}{#1}%
   #1\par}%
  \nopagebreak
  \vspace{\dimexpr5.65bp-\parskip\relax}}

% ---------------------------------------------------------------------
%  Diagramas (Excalidraw)
% ---------------------------------------------------------------------
%  \diagrama{nombre}{Título de la figura}
%  Inserta diagramas/dist/nombre.png, generado a partir de
%  diagramas/nombre.excalidraw con npm run diagramas. Al editar el diagrama y
%  volver a ejecutar el script, el documento se actualiza solo.
\newcounter{figuradiagrama}
\newcommand\diagrama[2]{%
  \par\addvspace{6bp}%
  \refstepcounter{figuradiagrama}\label{fig:#1}%
  \begingroup\centering
  \includegraphics[width=\linewidth,height=0.55\textheight,keepaspectratio]%
    {diagramas/dist/#1.png}\par
  \nopagebreak\vspace{4bp}%
  {\onceSencillo\textbf{Figura \thefiguradiagrama.} #2\par}%
  \endgroup
  \addvspace{6bp}}

% ---------------------------------------------------------------------
%  Tablas (CSV)
% ---------------------------------------------------------------------
%  \tabla{nombre}{Título de la tabla}{anchos de columna separados por comas}
%  Inserta la tabla definida en tables/nombre.csv: separador ";", la primera
%  fila es el encabezado y el resto son los datos. Los anchos son una dimensión
%  por columna, p. ej. {1.4cm,3.4cm,2.3cm,6.8cm,1.5cm}. La tabla se numera
%  ("Tabla N.") y puede citarse con \ref{tab:nombre}. Evitar en las celdas los
%  caracteres especiales de LaTeX (% & _ \ $) y el punto y coma.
\newcounter{tablacsv}
\ExplSyntaxOn
\tl_new:N \l__tabla_preamble_tl
\seq_new:N \l__tabla_head_seq
\cs_new_protected:Npn \__tabla_headrow:
  {
    \seq_set_map:NNn \l__tabla_head_seq \g__csvsim_line_seq { \textbf{##1} }
    \seq_use:Nn \l__tabla_head_seq { & }
  }
\cs_new_protected:Npn \__tabla_read:nn #1#2
  {
    \csvreader
      [
        separator = semicolon,
        respect~sharp,
        head,
        column~names~detection = false,
        after~head =
          {
            \begin{longtable}{#1}
            \hline \rowcolor{gray!25} \__tabla_headrow: \\ \hline \endfirsthead
            \hline \rowcolor{gray!25} \__tabla_headrow: \\ \hline \endhead
            \hline \endfoot
          },
        late~after~line = { \\ \hline },
        late~after~last~line = { \end{longtable} },
      ]
      {tables/#2.csv}{}{\csvlinetotablerow}
  }
\cs_generate_variant:Nn \__tabla_read:nn { V }
\NewDocumentCommand\tabla{m m m}
  {
    \tl_set:Nn \l__tabla_preamble_tl { | }
    \clist_map_inline:nn {#3}
      { \tl_put_right:Nn \l__tabla_preamble_tl { >{\raggedright\arraybackslash}p{##1} | } }
    \par\addvspace{6bp}
    \refstepcounter{tablacsv}\label{tab:#1}
    {\onceSencillo\textbf{Tabla~\thetablacsv.}~#2\par}
    \vspace{4bp}
    \begingroup
      \setlength{\tabcolsep}{3pt}\renewcommand{\arraystretch}{1.15}
      \fontsize{10bp}{12bp}\selectfont\raggedright
      \__tabla_read:Vn \l__tabla_preamble_tl {#1}
    \endgroup
  }
\ExplSyntaxOff

% ---------------------------------------------------------------------
%  Índice (estilos "TDC 1" y "TDC 2" de Word)
% ---------------------------------------------------------------------
%  TDC 1: Arial 11 mayúsculas, 1,5, tabulación derecha con puntos a
%         17,79 cm (5,55 pt más allá del margen derecho).
%  TDC 2: Arial 11, sangría 0,5 cm, tabulación derecha a 17,09 cm.
\makeatletter
\renewcommand*\l@section[2]{%
  \par{\onceUnoCinco\leftskip=0pt\rightskip=0pt\parfillskip=0pt
       #1\nobreak\xleaders\hbox{.}\hfill\nobreak\hbox{#2}\kern-5.55bp\par}}
\renewcommand*\l@subsection[2]{%
  \par{\onceUnoCinco\leftskip=14.2bp\rightskip=0pt\parfillskip=0pt
       #1\nobreak\xleaders\hbox{.}\hfill\nobreak\hbox{#2}\kern14.15bp\par}}
\newcommand\indiceGeneral{\@starttoc{toc}}
\makeatother

% ---------------------------------------------------------------------
%  Referencias bibliográficas (Arial 11, interlineado 1,15)
%  Se numeran por orden de primera cita: \cite{clave} -> [n]
% ---------------------------------------------------------------------
\urlstyle{same}
\makeatletter
\g@addto@macro\UrlBreaks{\do\-\do\_\do\=}
\newenvironment{referencias}{%
  \list{\@biblabel{\@arabic\c@enumiv}}{%
    \settowidth\labelwidth{[99]}%
    \leftmargin\labelwidth \advance\leftmargin\labelsep
    \itemsep=4bp \parsep=0pt \topsep=0pt \partopsep=0pt
    \usecounter{enumiv}\let\p@enumiv\@empty
    \renewcommand\theenumiv{\@arabic\c@enumiv}}%
  \onceUnoQuince \sloppy}{\endlist}
\makeatother

% Listas numeradas de Word: número a 0,63 cm, texto a 1,27 cm
\setlist[enumerate]{label=\arabic*.,leftmargin=36bp,labelwidth=18bp,
                    labelsep=0pt,align=left,topsep=0pt,partopsep=0pt,
                    itemsep=0pt}

\begin{document}

% =====================================================================
%  Sección 1 — Portada
% =====================================================================
\seccion{82.35bp}{14.15bp}{87.19bp}{28.37bp}{11.25bp}
\pagestyle{portada}
\pagenumbering{alph}
\doceUnoCinco
% Logotipo anclado: 5,48 cm desde el margen izquierdo y 1,04 cm bajo el
% inicio del párrafo; 6,85 x 2,78 cm.
\noindent\rule[-2.54bp]{0pt}{13.79bp}%
\rlap{\hspace{155.25bp}\smash{\raisebox{-96.85bp}{%
  \includegraphics[width=194.05bp,height=78.65bp]{imagenes/logo.jpg}}}}\par
\begingroup\centering
\timesnr\vacio\vacio\vacio\vacio\vacio\normalfont
\textbf{ Facultad de Tecnologías Interactivas}\par
\vacio\vacio
{\fontsize{24bp}{41.4bp}\selectfont \textbf{Sistema de microtransacciones mediante plataformas de pago nacionales}\par}
\vacio\vacio
{\fontsize{14bp}{24.15bp}\selectfont
\textbf{INFORME TÉCNICO}\par
ASIGNATURA: PROYECTO INVESTIGACIÓN Y DESARROLLO IV\par
DISCIPLINA: PRÁCTICA PROFESIONAL\par}
{\timesnr\doceSencillo\vacio\vacio}
\vacio\vacio\vacio
\textbf{Autor:} Alain\par
\textbf{Tutor(es):} MsC. Yadira Ramírez Rodríguez\par
\vacio\vacio
\ La Habana, Octubre de 2026\par
\vspace{5.65bp}
Año del centenario del comandante en jefe Fidel Castro Ruz\par
\endgroup

% =====================================================================
%  Sección 2 — Resumen / Abstract (numeración romana desde i)
% =====================================================================
\seccion{82.35bp}{14.15bp}{87.19bp}{28.37bp}{11.25bp}
\pagestyle{informe}
\renewcommand\textoEncabezado{Resumen}
\pagenumbering{roman}
\doceUnoCinco
\espacioInicial{11.35bp}
\rule[-2.54bp]{0pt}{13.79bp}\textbf{Resumen}\par
% Word no suma el espacio posterior de un párrafo y el anterior del
% siguiente: usa el mayor de los dos (11,35 pt).
{\setlength{\parskip}{11.35bp}
Los desarrolladores de aplicaciones y sitios web en Cuba encuentran dificultades para incorporar cobros en línea mediante las plataformas de pago nacionales, Transfermóvil y EnZona: trámites de contratación prolongados, documentación limitada y la necesidad de realizar una integración diferente con cada una. Como consecuencia, muchas ventas se gestionan mediante aplicaciones de mensajería, con pagos a cuentas personales que reducen la confianza del cliente y dificultan el comercio digital.

El objetivo de este trabajo es desarrollar un sistema web que permita a los desarrolladores cobrar microtransacciones desde sus aplicaciones y sitios web mediante una única interfaz de programación, que gestione la operación con la plataforma de pago nacional y comunique el resultado a la aplicación. Para ello se sistematizan los fundamentos teóricos, se analizan soluciones homólogas, se especifican los requisitos y se diseña la arquitectura de la solución con un enfoque iterativo e incremental. Se espera que el sistema facilite la incorporación de pagos digitales en las aplicaciones nacionales y contribuya a la seguridad y la trazabilidad de las operaciones.

\textbf{Palabras clave: }microtransacciones; pasarelas de pago; interfaz de programación de aplicaciones; webhooks; comercio electrónico; Cuba.\par}

\newpage
\textbf{\textit{Abstract}}\par
\vspace{11.35bp}
\textit{Software developers in Cuba face difficulties in adding online payments to their applications and websites through the national payment platforms, Transfermóvil and EnZona: lengthy contracting procedures, limited documentation, and the need to build a separate integration for each platform. As a result, many sales are handled through messaging applications, with payments sent to personal accounts, which reduces customer trust and hinders digital commerce.}\par
\vspace{11.35bp}
\textit{The aim of this work is to develop a web system that allows developers to charge microtransactions from their applications and websites through a single application programming interface, which manages the operation with the national payment platform and reports the result to the application. To this end, the theoretical foundations are systematized, similar solutions are analyzed, the requirements are specified and the architecture of the solution is designed following an iterative and incremental approach. The system is expected to make it easier to add digital payments to national applications and to improve the security and traceability of operations.}\par
\vacio
\textbf{\textit{Keywords: }}\textit{microtransactions; payment gateways; application programming interface; webhooks; e-commerce; Cuba.}\par
\vacio
\vacio

% =====================================================================
%  Sección 3 — Índice
% =====================================================================
\seccion{85bp}{16.8bp}{87.54bp}{28.72bp}{11.25bp}
\renewcommand\textoEncabezado{índice}
\doceUnoCinco
\espacioInicial{11.35bp}
\rule[-2.54bp]{0pt}{13.79bp}\textbf{Índice}\par
\vspace{11.35bp}
\indiceGeneral
\vacio

% =====================================================================
%  Sección 4 — Introducción (numeración arábiga desde 1)
% =====================================================================
\seccion{85bp}{16.8bp}{87.54bp}{28.72bp}{10.31bp}
\renewcommand\textoEncabezado{Introducción}
\pagenumbering{arabic}
\onceUnoCinco
\setlength{\parskip}{5.7bp}

\tituloI*{Introducción}

El comercio electrónico y los pagos digitales han transformado la forma en que se comercializan los productos y servicios, en particular los de naturaleza digital. En este ámbito, las microtransacciones, es decir, los pagos de pequeño importe realizados en línea, se han consolidado como modelo de ingresos de las aplicaciones y los videojuegos: según Tomić~\cite{tomic}, cerca de la mitad de los ingresos de las grandes editoras de videojuegos procede de ellas.

En Cuba, el proceso de bancarización ha impulsado la incorporación de canales de pago electrónico en los comercios~\cite{granma23}. Dos plataformas nacionales sustentan estos pagos: Transfermóvil, que comenzó a funcionar en 2015~\cite{etecsa}, y EnZona, puesta en explotación en julio de 2019~\cite{cubadebate21}. La primera superó los tres millones de usuarios en 2022~\cite{etecsa} y la segunda los 800 mil en 2023~\cite{mesa23b}. A ello se suma que la telefonía móvil llega a 7,8 millones de cubanos y que la cobertura poblacional de datos móviles supera el 85 \%~\cite{mesa23a}.

Sin embargo, la existencia de pasarelas de pago no asegura que los desarrolladores puedan incorporarlas a sus productos. Para cobrar en línea deben contratar el servicio con cada plataforma mediante procedimientos propios~\cite{granma23,oncuba23}. A esto se suman, según la experiencia del autor que se detalla en el epígrafe 1.3, la documentación pública limitada, la escasa comunicación con la comunidad de desarrolladores y la restricción de las direcciones IP con acceso al servicio. Por ello, muchas ventas se cierran mediante aplicaciones de mensajería, con pagos a cuentas personales no vinculadas a una entidad comercial ni a un trabajo por cuenta propia, lo que reduce la confianza del cliente y consume el tiempo del vendedor.

\textbf{Problema científico:} ¿Cómo facilitar a los desarrolladores cubanos la incorporación de microtransacciones en sus aplicaciones y sitios web mediante las plataformas de pago nacionales?

\textbf{Objeto de estudio:} Los procesos de cobro de productos y servicios digitales en aplicaciones y sitios web desarrollados en Cuba.

\textbf{Campo de acción:} La integración y gestión de microtransacciones mediante plataformas de pago nacionales en aplicaciones y sitios web desarrollados en Cuba, con énfasis en la generación de enlaces de pago, la confirmación de las transacciones y la notificación de sus resultados a las aplicaciones.

\textbf{Objetivo general:} Desarrollar un sistema web que permita a los desarrolladores de software cubanos cobrar por sus productos y servicios desde sus propias aplicaciones o sitios web mediante las plataformas de pago nacionales, de forma segura y con una única integración.

\textbf{Objetivos específicos:}
\begin{enumerate}[parsep=5.7bp]
  \item Sistematizar los fundamentos teóricos y las soluciones homólogas relacionados con las microtransacciones y la integración de pasarelas de pago.
  \item Caracterizar las limitaciones técnicas y de acceso del contexto nacional para la integración de plataformas de pago.
  \item Especificar los requisitos funcionales y no funcionales del sistema.
  \item Seleccionar los servicios externos, las bibliotecas y las herramientas con los que se desarrollará el sistema.
  \item Diseñar la arquitectura del sistema.
  \item Construir una primera versión funcional con los requisitos de criticidad alta.
  \item Probar el sistema a nivel funcional y no funcional para evaluar su calidad y su seguridad.
\end{enumerate}

\textbf{Hipótesis:} El desarrollo de un sistema que integre plataformas de pago nacionales y proporcione mecanismos para generar enlaces de pago y notificar la confirmación de las transacciones puede facilitar la incorporación de microtransacciones en aplicaciones y sitios web cubanos.

\textbf{Preguntas científicas:}
\begin{enumerate}[parsep=5.7bp]
  \item ¿Qué limitaciones técnicas y de acceso presentan las plataformas de pago nacionales para su integración en aplicaciones y sitios web cubanos?
  \item ¿Qué requisitos funcionales y de seguridad debe cumplir un sistema de gestión de microtransacciones en el contexto cubano?
  \item ¿Cómo diseñar una arquitectura que permita generar enlaces de pago y comunicar a las aplicaciones el resultado de las transacciones?
  \item ¿En qué medida la solución propuesta facilita la integración de microtransacciones y cumple los requisitos definidos?
\end{enumerate}

\textbf{Métodos científicos:} Para el desarrollo de la investigación se emplearán métodos teóricos y empíricos. Entre los teóricos, el analítico-sintético, en el estudio de las plataformas de pago nacionales, las soluciones homólogas y sus mecanismos de integración, y la modelación, para representar la arquitectura, los datos y las interacciones de la solución. Entre los empíricos, la observación, en la caracterización de los procesos de cobro actuales; la entrevista semiestructurada a desarrolladores y responsables de emprendimientos, para conocer cómo gestionan los cobros y qué dificultades encuentran; la revisión documental de fuentes oficiales y técnicas; y las pruebas de software, para evaluar el cumplimiento de los requisitos definidos, cuyos resultados se registrarán mediante una matriz de casos de prueba.

\textbf{Aporte práctico:} Un sistema que permita a los desarrolladores de software cubanos cobrar por su trabajo de forma sencilla, legal y transparente puede contribuir al crecimiento del mercado de aplicaciones nacionales y ofrecer una alternativa nacional a servicios que hoy se contratan a empresas de otros países.

\textbf{Estructura del documento:} El informe se organiza en tres capítulos. El Capítulo I presenta los fundamentos teóricos, el análisis de soluciones homólogas, el diagnóstico de la situación actual y el entorno de desarrollo. El Capítulo II expone las necesidades del cliente, la propuesta de solución técnica, los requisitos y el diseño de la arquitectura. El Capítulo III aborda el diseño de la base de datos, los prototipos funcionales de interfaz y la codificación.

{\fontsize{11bp}{13.65bp}\selectfont\vacio\vacio}

% =====================================================================
%  Sección 5 — Capítulo I
% =====================================================================
\seccion{70.75bp}{2.55bp}{87.54bp}{28.72bp}{11.25bp}
\renewcommand\textoEncabezado{Capítulo I}
\doceUnoCinco
\setlength{\parskip}{0pt}

\tituloI{Capítulo I: Marco teórico y diagnóstico}

En este capítulo se presentan los fundamentos teóricos que sustentan el desarrollo de un sistema de microtransacciones mediante plataformas de pago nacionales. Se abordan los conceptos relacionados con las microtransacciones, las interfaces de programación, las pasarelas de pago, las notificaciones automáticas y el contexto cubano de los pagos electrónicos. Asimismo, se analizan soluciones similares para identificar sus aportes y limitaciones respecto al contexto cubano. Posteriormente, se describe el contexto en el que se desarrolla la investigación y las dificultades existentes para incorporar pagos en aplicaciones y sitios web nacionales. Finalmente, se exponen la metodología y el entorno de desarrollo seleccionados para construir la propuesta.

\vacio
\tituloII{1.1 Fundamentos y referentes teórico-metodológicos sobre el objeto de estudio.}

Para fundamentar la propuesta es necesario precisar, en primer lugar, qué se entiende por microtransacción y qué características la distinguen, pues de ellas depende lo que el sistema debe soportar. A continuación se describen los elementos tecnológicos que permiten cobrar una microtransacción desde una aplicación: la interfaz de programación mediante la cual esta solicita el cobro, la pasarela de pago que lo procesa y el mecanismo de notificación que comunica el resultado. Por último, se caracteriza el contexto cubano de los pagos electrónicos en el que dichos elementos deben integrarse.

\vacio
\textbf{1.1.1 Microtransacciones}

\textit{Concepto.} Una microtransacción es un micropago, es decir, un pago de una pequeña cantidad de dinero realizado en línea, típicamente en un entorno de juego~\cite{rita24}. Tomić~\cite{tomic} la define como el pago que se realiza por adquirir una aplicación móvil o contenido adicional dentro de un videojuego. En la bibliografía el término se asocia sobre todo con los videojuegos y las aplicaciones móviles, aunque sus características son extensibles a cualquier producto o servicio digital de bajo valor.

\textit{Características.} De la bibliografía consultada se desprenden las siguientes:
\begin{itemize}[parsep=5.7bp]
  \item \textbf{Importe reducido y compra repetible.} Los pagos son de poca cuantía y se realizan de forma sucesiva, a medida que el usuario adquiere contenido adicional~\cite{tomic,rita24}.
  \item \textbf{Bien digital ligado a su plataforma.} Lo que se adquiere es un bien virtual cuyo valor depende del entorno en el que puede utilizarse~\cite{hamarik17}. Por ello, el pago debe confirmarse con rapidez para que la aplicación entregue el bien de forma inmediata.
  \item \textbf{Sensibilidad a la fricción del pago.} Odlyzko~\cite{odlyzko03} sostiene que los micropagos enfrentan obstáculos económicos, sociológicos y psicológicos, y no solo tecnológicos, que limitan su adopción. En consecuencia, un sistema de microtransacciones debe reducir al mínimo los pasos que el cliente realiza para pagar.
  \item \textbf{Costo del servicio de cobro.} En el contexto nacional, el costo de los servicios de pago electrónico depende del plan contratado~\cite{granma23}, lo que puede incidir en la viabilidad de cobrar importes pequeños.
\end{itemize}

\textit{Tipos.} Zendle, Meyer y Ballou~\cite{zendle20} distinguen las microtransacciones cosméticas, que modifican la apariencia sin ofrecer ventaja alguna, de las de tipo \textit{pay to win}, que aumentan las posibilidades de éxito del usuario; además estudian las cajas de botín (\textit{loot boxes}), cuyo contenido no se conoce de antemano. Con frecuencia se emplea una moneda virtual como intermediaria, una de cuyas formas, la moneda \textit{premium}, solo puede obtenerse mediante pago~\cite{tomic}.

\textit{Modelos de negocio y motivaciones.} Las microtransacciones sustentan el modelo \textit{free-to-play}, en el que el producto se obtiene sin costo y los ingresos provienen de la venta de bienes virtuales~\cite{hamari15,tomic}. Hamari~\cite{hamari15} halló que la actitud hacia los bienes virtuales y la opinión del entorno social aumentan la intención de compra, mientras que el disfrute del juego la reduce; y la revisión de Hamari y Keronen~\cite{hamarik17} confirma que las razones de compra están estrechamente ligadas a la plataforma donde se venden los bienes.

En la presente investigación, las microtransacciones constituyen el tipo de operación que gestionará el sistema propuesto. Este facilitará que los desarrolladores reciban pagos mediante plataformas nacionales y que sus aplicaciones conozcan el resultado de cada operación para entregar el producto o servicio correspondiente. Para cobrar estas operaciones, la aplicación debe poder solicitar el pago a otro sistema, lo que conduce al concepto de interfaz de programación.

\vacio
\textbf{1.1.2 Interfaces de programación de aplicaciones (API)}

Una interfaz de programación de aplicaciones (API, por sus siglas en inglés) define un conjunto de operaciones y reglas que permite la comunicación entre sistemas de software, sin que uno deba conocer los detalles de implementación interna del otro. En la web, las API suelen apoyarse en el estilo arquitectónico REST propuesto por Fielding~\cite{fielding00}, que se basa en la identificación de recursos accesibles mediante HTTP y establece restricciones como la separación entre cliente y servidor y la comunicación sin estado entre solicitudes. Los proveedores de pagos exponen sus servicios mediante interfaces de este tipo; por ejemplo, Stripe permite crear enlaces de pago que el comercio comparte con sus clientes~\cite{stripelinks}.

En la presente investigación, las API permitirán la comunicación entre las aplicaciones de los desarrolladores y el sistema propuesto, así como la integración de este con las plataformas de pago nacionales. A través de estas interfaces se podrán solicitar enlaces de pago y consultar el estado de las operaciones, de acuerdo con las capacidades que ofrezca cada plataforma. La solicitud recibida por la API debe encaminarse hacia el sistema financiero, función que cumplen las pasarelas de pago.

\textbf{1.1.3 Pasarelas de pago}

Según el Banco de la Reserva de la India~\cite{rbi19}, una pasarela de pago es un proveedor de infraestructura tecnológica que encamina y facilita el procesamiento de una transacción de pago en línea, sin intervenir en el manejo de los fondos. Esa misma fuente señala que las pasarelas y los agregadores de pago permiten a los comercios aceptar distintos instrumentos de pago sin necesidad de crear un sistema de integración propio. Esta definición distingue la pasarela de las entidades bancarias que ejecutan los movimientos de dinero.

En la presente investigación, las plataformas de pago nacionales constituyen los servicios externos con los que se busca integrar el sistema propuesto. Este actuará como una capa de integración para facilitar que los desarrolladores incorporen cobros en sus aplicaciones, sin sustituir las funciones de las plataformas ni las de las entidades bancarias. El alcance de dicha integración dependerá de las interfaces disponibles y de las condiciones de acceso de cada proveedor. Una vez procesado el pago, la aplicación debe conocer su resultado, para lo cual se emplean las notificaciones automáticas.

\textbf{1.1.4 Notificaciones automáticas mediante webhooks}

Un webhook es un mecanismo de comunicación mediante el cual un sistema envía una solicitud HTTP a una dirección previamente configurada cuando ocurre un evento determinado. Stripe lo describe como un punto de acceso HTTPS al que la plataforma envía, en formato JSON, los eventos ocurridos en la cuenta del comercio, lo que permite responder a eventos asíncronos, como la confirmación bancaria de un pago~\cite{stripeweb}. De este modo, el sistema receptor no tiene que consultar continuamente si se han producido novedades.

Para que estas notificaciones puedan utilizarse de forma confiable es necesario verificar su autenticidad, para lo cual se emplean firmas basadas en códigos de autenticación de mensajes con clave (HMAC)~\cite{rfc2104}, y contemplar los reintentos mediante un procesamiento idempotente, es decir, evitar que la recepción repetida de una misma notificación provoque la ejecución duplicada de una operación. Stripe aplica el mismo principio a las solicitudes de su API, que pueden reintentarse sin ejecutar dos veces la misma operación~\cite{stripeidem}. Además, la criticidad de estas interfaces exige controles de autenticación y autorización; el proyecto OWASP~\cite{owasp23} recoge entre los principales riesgos de seguridad de las API la autenticación defectuosa y la autorización deficiente a nivel de objeto.

En la presente investigación, los webhooks permitirán que el sistema propuesto notifique a las aplicaciones de los desarrolladores los cambios relevantes en el estado de las transacciones. La confirmación del pago deberá sustentarse en información validada de la plataforma correspondiente. Los elementos descritos deben operar en el contexto cubano, que se caracteriza a continuación.

\vacio
\textbf{1.1.5 Los pagos electrónicos en el contexto cubano}

Una de las medidas del Banco Central de Cuba en el camino hacia la bancarización es que los comercios que venden bienes y servicios incorporen canales de pago electrónico~\cite{granma23}. Dos plataformas sustentan estos pagos. Transfermóvil, desarrollada por ETECSA, comenzó a funcionar en 2015~\cite{etecsa}; requiere una tarjeta matriz de los bancos Metropolitano, de Crédito y Comercio o Popular de Ahorro, y permite operar sin activar los datos móviles~\cite{etecsat}. EnZona, desarrollada por XETID, se puso en explotación en julio de 2019 y en octubre de ese año incorporó el módulo de comercio electrónico con pago mediante código QR~\cite{cubadebate21}; funciona con tarjetas de los mismos bancos, necesita conexión de datos, aunque su acceso es libre de consumo en la red nacional~\cite{xetid,mesa23b}.

Para integrarlas, los desarrolladores deben contratar el servicio. En Transfermóvil, la modalidad de tienda virtual permite el cobro en línea y la integración a su API, y la contratación se realiza de forma digital~\cite{oncuba23}. En EnZona, el comercio se registra en el Bulevar mediante cinco pasos y se suscribe a un plan, entre los que se encuentran la pasarela de pagos y los códigos QR estáticos y dinámicos, cuyo costo se calcula según la modalidad elegida~\cite{granma23}; además, XETID ofrece una interfaz para desarrollar aplicaciones con las API de la plataforma~\cite{xetid}. La penetración de la telefonía móvil, con 7,8 millones de cubanos con teléfono, y la cobertura de datos móviles, superior al 85 \% de la población~\cite{mesa23a}, ofrecen una base para estos servicios.

Se observa así que existen plataformas y vías de integración, pero que cada una exige su propio proceso de contratación y su propia integración. Conocidos los fundamentos y el contexto, resulta necesario estudiar cómo se resuelve en la práctica el cobro en línea, lo que se aborda en el epígrafe siguiente.

\vacio
\tituloII{1.2 Soluciones homólogas.}

El análisis de soluciones homólogas permite identificar experiencias y mecanismos de integración que pueden servir como referentes para el desarrollo del sistema propuesto. En esta investigación se consideran tres: Stripe, como referente internacional de integración de pagos para desarrolladores, y Transfermóvil y EnZona, por ser las plataformas nacionales con las que se integrará el sistema. Aunque presentan diferencias en su alcance y funcionamiento, se comparan según la forma de integración, la generación de enlaces de pago, la notificación de resultados y las condiciones de acceso.

\vacio
\textbf{1.2.1 Descripción y análisis de las particularidades relevantes}

\textbf{Stripe}

Stripe es una plataforma de servicios de pago que ofrece herramientas para incorporar cobros en aplicaciones y sitios web. Permite crear enlaces de pago que el comercio comparte en redes sociales, correos electrónicos o su sitio web, sin escribir código~\cite{stripelinks}; envía los eventos ocurridos a un punto de acceso HTTPS del comercio mediante webhooks~\cite{stripeweb}; y admite solicitudes idempotentes en su API~\cite{stripeidem}.

Su principal interés para esta investigación radica en la experiencia de integración que proporciona a los desarrolladores y en sus mecanismos para gestionar el ciclo de vida de los pagos. Constituye un referente para estudiar la documentación de una API, la comunicación de cambios de estado y el tratamiento de solicitudes repetidas. Se toma como referente funcional del diseño de la API, sin presuponer su disponibilidad para desarrolladores establecidos en Cuba.

\vacio
\textbf{Transfermóvil}

Transfermóvil es una solución nacional desarrollada por ETECSA que permite realizar pagos y otras operaciones desde dispositivos móviles, entre ellas el pago en tiendas virtuales~\cite{etecsat}. En su modalidad de tienda virtual, admite el cobro en línea y la integración a su API; el interesado contrata el servicio de forma digital y, una vez validada la solicitud, ETECSA o el Grupo de Electrónica para el Turismo habilitan el servicio~\cite{oncuba23}. Es, además, la plataforma con mayor número de usuarios~\cite{etecsa,mesa23b}.

En esta investigación interesa estudiar sus mecanismos de pago a comercios y las condiciones técnicas y contractuales para incorporar dichos mecanismos en aplicaciones de terceros. Debe distinguirse entre las operaciones que puede realizar un usuario desde la aplicación y los servicios de integración disponibles para un desarrollador. Las posibilidades de automatizar la confirmación de pagos deberán determinarse a partir de documentación oficial y de las condiciones de acceso al servicio.

\vacio
\textbf{EnZona}

EnZona es una plataforma cubana de comercio electrónico, desarrollada por XETID, que permite realizar y recibir transferencias y pagos con tarjetas de los bancos Metropolitano, BANDEC y BPA~\cite{xetid}. Para los comercios ofrece servicios de pasarela de pagos, códigos QR estáticos y dinámicos y Bulevar virtual~\cite{granma23}, y XETID declara una interfaz para desarrollar aplicaciones con sus API~\cite{xetid}. Según testimonios recogidos por la prensa, resulta más fácil de integrar con sitios web de terceros que otras alternativas nacionales~\cite{eltoque20}.

Para esta investigación resulta relevante examinar los mecanismos de autenticación, creación de operaciones y consulta o notificación de sus resultados que se encuentren disponibles en su documentación. También deben analizarse los requisitos de contratación y las condiciones de acceso, pues estos elementos delimitan las posibilidades reales de integración con el sistema propuesto.

\vacio
\textbf{Síntesis del análisis}

La Tabla~\ref{tab:comparacion-homologas} resume la comparación. Cuando las fuentes consultadas no permiten afirmar que una característica esté disponible, se indica que debe confirmarse con la documentación oficial.

\vacio
\tabla{comparacion-homologas}{Comparación de las soluciones homólogas.}{2.8cm,4.2cm,4.2cm,4.2cm}

\vacio
A partir de la comparación se determina qué se toma de cada solución y qué no. De Stripe se toma su modelo de integración: enlaces de pago, notificaciones mediante webhooks con verificación de su autenticidad y solicitudes idempotentes. No se trasladan sus medios de pago, pues el sistema opera con las plataformas y las tarjetas de los bancos nacionales. De Transfermóvil y EnZona se toman las vías de integración y las condiciones de contratación, que delimitan el alcance real del sistema. No se parte de que ofrezcan enlaces de pago ni notificaciones equivalentes a las de Stripe, por lo que el sistema deberá generar sus propios enlaces y confirmar el resultado de cada operación con la plataforma correspondiente.

Con ello, el sistema propuesto se concibe como una capa de integración orientada a simplificar la incorporación de pagos nacionales en aplicaciones y sitios web. Su alcance no consiste en sustituir las plataformas existentes, sino en facilitar su utilización mediante una interfaz común, la gestión de enlaces de pago y la notificación del resultado de las operaciones, siempre que las capacidades y condiciones de cada proveedor lo permitan. Estos resultados se contrastan a continuación con un caso concreto.

\vacio
\tituloII{1.3 Diagnóstico organizacional.}

La investigación se desarrolla en el contexto de aplicaciones como “El Yerro Menú”, cuyos productos y servicios se comercializan mediante su sitio web. Para caracterizar la gestión actual de los cobros se empleó la observación del proceso, con el propósito de identificar los actores involucrados, las actividades del proceso y las dificultades existentes.

\vacio
Actualmente, el proceso de cobro consiste en ponerse en contacto con el proveedor para que este le explique las condiciones del pago, y los métodos que puede utilizar, usualmente enviando un número de tarjeta y un número de teléfono para confirmar el pago. A partir de la información recopilada se identificó que el proveedor pierde demasiado tiempo en interactuar con clientes que no están realmente interesados en realizar la compra, además de que la confianza del cliente se ve afectada al momento de enviar una transferencia bancaria a una cuenta bancaria de un desconocido.

\vacio
Por otra parte, a partir de su experiencia como desarrollador, el autor identificó dificultades para contratar y utilizar las API de pago de las dos plataformas nacionales.

\textbf{Transfermóvil.} Es la plataforma con mayor número de usuarios, por lo que resulta la más habitual entre los clientes. Sin embargo, la contratación de su API pasa por varios inconvenientes:
\begin{enumerate}[parsep=5.7bp]
  \item Falta de documentación pública.
  \item Falta de comunicación con la comunidad de desarrolladores.
  \item Sistema complejo de implementación.
  \item Regulación de las direcciones IP que pueden tener acceso al servicio.
\end{enumerate}

\textbf{EnZona.} Su adopción por parte de la población es menor, pero es la plataforma que mejor acceso brinda a los desarrolladores a su API. Esto no significa que sea perfecta, pues presenta inconvenientes como:
\begin{enumerate}[parsep=5.7bp]
  \item Documentación pública existente, pero en algunos casos deficiente.
  \item Falta de comunicación con la comunidad de desarrolladores.
  \item Regulación de las direcciones IP que pueden tener acceso al servicio.
\end{enumerate}

\vacio
Estas insuficiencias evidencian la necesidad de facilitar la integración con las plataformas de pago nacionales y automatizar, cuando sus capacidades lo permitan, la confirmación de las operaciones y la comunicación de sus resultados. Para construir la solución es necesario definir ahora cómo se desarrollará.

\vacio
\tituloII{1.4 Metodología y entorno de desarrollo.}

El desarrollo del sistema se organizará mediante un enfoque iterativo e incremental, cuyo rasgo común, según Larman y Basili~\cite{larman03}, es evitar un proceso secuencial de una sola pasada, dirigido por documentos y por etapas cerradas. En cada iteración se analizan, diseñan, implementan y prueban un conjunto reducido de requisitos priorizados, de modo que el sistema crece de forma gradual y puede ajustarse a las condiciones técnicas de los servicios externos. Las primeras iteraciones priorizan los mecanismos de integración con las plataformas de pago nacionales, la generación de enlaces de pago y la notificación del resultado de las operaciones.

Dada la criticidad de un sistema que gestiona operaciones de pago, el entorno de desarrollo se compone de tecnologías seguras, eficientes y de probada madurez:

\begin{itemize}[parsep=5.7bp]
  \item \textbf{TypeScript:} lenguaje que extiende JavaScript con un sistema de tipos~\cite{typescript}, lo que permite detectar errores durante el desarrollo y emplear un único lenguaje en la interfaz de usuario y en la API.
  \item \textbf{Next.js:} marco de trabajo para aplicaciones web basado en React~\cite{nextjs}, que permite construir en un mismo proyecto la interfaz de usuario y la API del sistema, lo que simplifica el desarrollo, el despliegue y el mantenimiento.
  \item \textbf{PostgreSQL:} sistema de gestión de bases de datos relacional de código abierto. Sus transacciones agrupan varios pasos en una única operación de todo o nada~\cite{postgres}, característica fundamental para el registro de pagos.
  \item \textbf{Prisma ORM:} herramienta de mapeo objeto-relacional que facilita el acceso a la base de datos mediante consultas tipadas y gestiona la evolución del esquema mediante migraciones~\cite{prisma}.
  \item \textbf{Zod:} biblioteca de validación de esquemas para TypeScript~\cite{zod}, que verifica la estructura y el contenido de los datos recibidos por la API y por los webhooks antes de procesarlos, lo que contribuye a la seguridad del sistema.
  \item \textbf{MinIO:} servicio de almacenamiento de objetos compatible con la API de Amazon S3~\cite{minio}, que se empleará para guardar los archivos asociados al sistema, como las imágenes de los productos y servicios, y que puede desplegarse en infraestructura propia.
  \item \textbf{Git:} sistema de control de versiones distribuido~\cite{chacon14}, que permite registrar la evolución del código en cada iteración.
  \item \textbf{Excalidraw:} herramienta para elaborar los diagramas que modelan la arquitectura de la solución.
\end{itemize}

\vacio
Estas herramientas se complementan entre sí, cuentan con documentación pública y pueden desplegarse en infraestructura propia, lo que resulta pertinente dado el contexto nacional en que deberá operar la solución.

\vacio
\tituloII{Conclusiones del capítulo.}

El estudio de los fundamentos teóricos permite reconocer las características de las microtransacciones y la importancia de las API, las pasarelas de pago y las notificaciones automáticas para gestionarlas en aplicaciones y sitios web. Su articulación constituye una base para diseñar una solución que comunique el resultado de los pagos y facilite la entrega de los productos o servicios correspondientes.

El análisis de las soluciones homólogas proporciona referentes para la integración y gestión de pagos. No obstante, las características de la solución internacional no pueden trasladarse directamente al contexto cubano sin considerar las capacidades y condiciones de acceso de las plataformas nacionales.

El diagnóstico evidencia las dificultades de cobro que enfrentan los desarrolladores y los clientes en el contexto estudiado, aunque será necesario complementarlo con evidencias adicionales que permitan precisar sus necesidades.

Finalmente, el enfoque iterativo e incremental y el entorno de desarrollo seleccionado permitirán relacionar el estudio del problema con el diseño y la evaluación del sistema. Estos elementos orientarán la definición de las necesidades, los requisitos y la arquitectura que se abordarán en el capítulo siguiente.

% =====================================================================
%  Sección 6 — Capítulo II
% =====================================================================
\seccion{73.6bp}{5.4bp}{87.54bp}{28.72bp}{11.25bp}
\renewcommand\textoEncabezado{Capítulo II}
\doceUnoCinco

\tituloI{Capítulo II: Necesidades del cliente, propuesta de solución técnica}

En este capítulo se presenta el diseño de la solución a partir de los referentes teóricos, las soluciones homólogas y el diagnóstico organizacional expuestos en el Capítulo I. Primero se identifican y analizan las necesidades del cliente, tomando como punto de partida las dificultades que enfrentan los desarrolladores cubanos para cobrar por sus productos y servicios. Después se describe la propuesta de solución técnica, que adapta al contexto nacional las buenas prácticas de las soluciones homólogas, en particular la integración mediante una interfaz común, la generación de enlaces de pago y la notificación automática de los resultados. Seguidamente se especifican los requisitos funcionales y no funcionales que debe cumplir el sistema y, por último, se expone el diseño de la arquitectura que los sustenta, junto con las decisiones adoptadas para definirla.

\vacio
\tituloII{2.1 Necesidades del cliente.}

% Este epígrafe se dedica a la identificación, descripción y análisis de las necesidades del cliente, precisando sus expectativas, restricciones y prioridades, así como su relación con el problema de investigación y con las insuficiencias identificadas en el diagnóstico organizacional.

Para cobrar en línea, los desarrolladores cubanos de aplicaciones y sistemas web deben solicitar el acceso a las pasarelas de pago nacionales mediante los formularios de registro oficiales. En ocasiones, la respuesta a estas solicitudes tarda semanas debido a trámites burocráticos, lo que retrasa el inicio de la comercialización de sus productos. Además, durante la implementación deben adaptar su sistema al funcionamiento particular de cada una de las dos plataformas, lo que incrementa el esfuerzo y los tiempos de desarrollo.

\vacio

En consecuencia, el cliente necesita un sistema que le permita cobrar por sus productos y servicios de forma ágil, con una única integración y sin depender de largos plazos de espera ni de implementaciones complejas.

\vacio
\tituloII{2.2 Propuesta de solución técnica.}

%Este epígrafe se dedica a la descripción de la propuesta de solución técnica, precisando su alcance, componentes, tecnologías seleccionadas y relación con las necesidades del cliente y con el problema de investigación.

Para responder a la necesidad descrita en el epígrafe anterior, se propone el desarrollo de un sistema web que permita a los desarrolladores de software cubanos incorporar microtransacciones en sus aplicaciones y sitios web sin tener que integrarse por separado con cada plataforma de pago nacional. El sistema actúa como una capa intermedia: la aplicación del desarrollador se comunica únicamente con una interfaz común, y el sistema se encarga de gestionar la operación con la plataforma de pago correspondiente y de informar el resultado a la aplicación.

\vacio
Mediante el sistema, el desarrollador podrá:
\begin{itemize}[parsep=5.7bp]
  \item Darse de alta y gestionar su cuenta.
  \item Registrar sus aplicaciones y obtener las credenciales para comunicarse con la API.
  \item Registrar la cuenta bancaria del trabajador por cuenta propia (TCP) o de la empresa que recibirá los pagos.
  \item Registrar sus productos y servicios, con su descripción y precio.
  \item Consultar el historial de los pagos realizados por cada producto o servicio.
  \item Configurar los webhooks mediante los cuales sus aplicaciones serán notificadas de las compras realizadas.
\end{itemize}

\vacio  
El flujo de cobro propuesto se desarrolla en cuatro pasos. Primero, la aplicación del desarrollador solicita al sistema, mediante su API, un enlace de pago para un producto o servicio previamente registrado. Segundo, el cliente final accede a dicho enlace y realiza el pago a través de la plataforma nacional correspondiente. Tercero, el sistema confirma el resultado de la operación a partir de la información validada de la plataforma y actualiza su estado. Por último, el sistema notifica el resultado a la aplicación mediante un webhook, de modo que esta pueda entregar el producto o servicio. Con este flujo se eliminan la comunicación manual con el proveedor y el envío de datos bancarios a un desconocido, dificultades identificadas en el diagnóstico organizacional.

\vacio
Dada la criticidad de un sistema que gestiona operaciones de pago, su desarrollo se apoya en las tecnologías seleccionadas en el epígrafe 1.4: Next.js para la interfaz de usuario y la API, PostgreSQL con Prisma ORM para los datos, Zod para validar la información recibida y MinIO para almacenar los archivos, todo ello en el lenguaje TypeScript. Su uso conjunto en un único lenguaje reduce la complejidad del desarrollo. La solución propuesta responde así a la necesidad de los desarrolladores de cobrar por sus productos y servicios de forma ágil, y constituye la base para la especificación de requisitos y el diseño de la arquitectura que se presentan a continuación.

\vacio
\tituloII{2.3 Especificación formal de requisitos funcionales y no funcionales.}

Este epígrafe consta de la especificación formal de los requisitos funcionales y no funcionales resultantes de la Ingeniería de Requisitos desarrollada, incluyendo su identificación, descripción, priorización y trazabilidad con las necesidades del cliente y con la propuesta de solución técnica.

\vacio
\textbf{2.3.1 Requisitos funcionales}

A partir de las necesidades del cliente y de la propuesta de solución técnica, se identificaron los requisitos funcionales que debe cumplir el sistema. Se describen en la Tabla~\ref{tab:requisitos-funcionales} y se organizan según los actores que interactúan con él: el desarrollador, el administrador, el cliente final, la aplicación del desarrollador mediante la API y el propio sistema. La prioridad (Alta, Media o Baja) refleja la importancia del requisito para la primera versión del software.

\vacio
\tabla{requisitos-funcionales}{Requisitos funcionales del sistema.}{1.35cm,3.4cm,2.3cm,6.8cm,1.5cm}

\vacio
\tituloII{2.4 Diseño de la arquitectura.}

Este epígrafe presenta el diseño de la arquitectura de la solución propuesta, incluyendo la descripción de los componentes, sus interrelaciones, los patrones de diseño empleados y los diagramas correspondientes, así como la justificación de las decisiones arquitectónicas adoptadas.

\vacio
Para estructurar el sistema se adopta el patrón arquitectónico Modelo-Vista-Controlador (MVC), que separa la interfaz de usuario, la lógica de control y la lógica de negocio con su acceso a los datos. Esta separación facilita el mantenimiento y las pruebas del sistema, y permite que cambios en una capa afecten lo mínimo posible a las demás. La Figura~\ref{fig:mvc-paquetes} muestra el diagrama de paquetes resultante.

\vacio
\diagrama{mvc-paquetes}{Diagrama de paquetes de la arquitectura MVC del sistema.}

\vacio
\tituloII{Conclusiones del capítulo.}

En este capítulo se definieron las bases técnicas de la solución a partir del diagnóstico realizado en el Capítulo I. Del análisis efectuado se concluye lo siguiente:

\begin{itemize}[parsep=5.7bp]
  \item La identificación de las necesidades del cliente evidenció que los desarrolladores cubanos requieren cobrar por sus productos y servicios de forma ágil, con una única integración y sin depender de largos plazos de espera ni de adaptaciones complejas a cada plataforma de pago nacional.
  \item La propuesta de solución técnica, un sistema web que actúa como capa intermedia entre las aplicaciones de los desarrolladores y las plataformas de pago, da respuesta a dichas necesidades mediante una interfaz común, la generación de enlaces de pago y la notificación automática de los resultados. Su desarrollo se apoya en tecnologías maduras y compatibles entre sí (Next.js, PostgreSQL, Prisma ORM, Zod y MinIO) que comparten un único lenguaje, TypeScript.
  \item La especificación de los requisitos funcionales, organizados por actores y priorizados según su importancia para la primera versión, permite delimitar el alcance del sistema y mantener la trazabilidad con las necesidades del cliente.
  \item El diseño de la arquitectura, basado en el patrón Modelo-Vista-Controlador, separa las responsabilidades del sistema en capas, lo que facilita su mantenimiento y sus pruebas, y constituye el fundamento para el diseño de la base de datos, los prototipos de interfaz y la codificación que se abordan en el capítulo siguiente.
\end{itemize}

\vacio
\vacio
\vacio

% =====================================================================
%  Sección 7 — Capítulo III
% =====================================================================
\seccion{70.75bp}{2.55bp}{87.54bp}{28.72bp}{11.25bp}
\renewcommand\textoEncabezado{Capítulo III}
\doceUnoCinco

\tituloI{Capítulo III (opcional): Diseño de base de datos, prototipos funcionales de interfaz, codificación}

El presente capítulo tiene como objetivo presentar el diseño de la base de datos, los prototipos funcionales de interfaz y la codificación de la solución propuesta. A lo largo de su desarrollo se abordan los principales contenidos siguientes: el diseño de la base de datos que sustenta el almacenamiento y la integridad de la información; el diseño de los prototipos funcionales de interfaz que materializan la interacción con el usuario; y la codificación de la solución, incluyendo los componentes, los estándares de codificación y el código generado. El capítulo se estructura en epígrafes que permiten al lector transitar desde el diseño de los datos y las interfaces hasta la implementación concreta de la solución, culminando con las conclusiones que sintetizan la pertinencia de las decisiones adoptadas. Como texto introductorio, se anticipa que estos elementos son consecuencia directa de las necesidades del cliente, la propuesta de solución técnica, la especificación formal de requisitos y el diseño de la arquitectura desarrollados en el Capítulo II.

\vacio
\tituloII{3.1 Diseño de la base de datos.}

Este epígrafe presenta el diseño de la base de datos de la solución propuesta, incluyendo el modelo de datos, las entidades y relaciones, las restricciones de integridad, la normalización, los mecanismos de almacenamiento y seguridad, así como los diagramas correspondientes y la justificación de las decisiones adoptadas.

\vacio
\tituloII{3.2 Prototipos funcionales de interfaz.}

Este epígrafe presenta el diseño de los prototipos funcionales de interfaz de la solución propuesta, incluyendo la descripción de las pantallas, la navegación, la usabilidad y la experiencia de usuario, así como los diagramas y esquemas correspondientes que materializan la interacción con el usuario.

\vacio
\tituloII{3.3 Codificación.}

Este epígrafe presenta el diseño de los componentes y la descripción de los estándares de codificación empleados. Se representa el código generado y la solución obtenida, precisando la correspondencia entre los componentes implementados, los requisitos especificados y la arquitectura diseñada.

\vacio
\tituloII{Conclusiones del capítulo.}

Las conclusiones se orientan a establecer los argumentos que respaldan la posición científica que adopta el autor a partir de la aplicación del enfoque, metodología y métodos de Ingeniería de Software aplicados para lograr cumplir el objetivo, resolver el problema científico planteado y obtener un resultado práctico, destacando la pertinencia de las decisiones tomadas en el diseño de la base de datos, los prototipos funcionales de interfaz y la codificación de la solución propuesta.

\vacio
\vacio
\vacio

% =====================================================================
%  Sección 8 — Conclusiones
%  (el margen superior de Word, 2,40 cm, queda por debajo del encabezado;
%   el texto empieza justo después de este, como hace Word)
% =====================================================================
\seccion{68.2bp}{0bp}{87.54bp}{28.72bp}{11.25bp}
\renewcommand\textoEncabezado{Conclusiones finales}
\doceUnoCinco

\tituloI{Conclusiones}

<Las conclusiones, por lo general van dirigidas a establecer los argumentos y resultados a los que se arribó en los siguientes aspectos: (1) sistematización del estado del arte referido al objeto de estudio y el campo de acción, (2) diagnóstico del estado actual del objeto de estudio, (3) principales aspectos del análisis, diseño e implementación de la solución, (4) principales resultados de la validación de la solución propuesta. Deben apoyarse en los resultados obtenidos y descritos en la memoria y no en datos que no aparezcan en este documento. No pueden exceder una cuartilla en su extensión>

\vacio
\vacio
\vacio

% =====================================================================
%  Sección 9 — Recomendaciones
% =====================================================================
\seccion{68.2bp}{0bp}{87.54bp}{28.72bp}{11.25bp}
\renewcommand\textoEncabezado{Recomendaciones}
\doceUnoCinco

\tituloI{Recomendaciones}

<Las recomendaciones, por lo general, van dirigidas a establecer aquellos aspectos en los cuales la investigación puede continuar para su perfeccionamiento, mantenimiento o evolución en el tiempo. No deben constituir acciones no realizadas por omisión de etapas del proceso investigativo o ingenieril; ni ser demasiadas en número que cuestionen la completitud y pertinencia de la investigación realizada. No pueden exceder una cuartilla en su extensión>

\vacio
\vacio
\vacio

% =====================================================================
%  Sección 10 — Referencias bibliográficas (Arial 11)
% =====================================================================
\seccion{85bp}{16.8bp}{87.54bp}{28.72bp}{10.31bp}
\renewcommand\textoEncabezado{Referencias bibliográficas}
\doceUnoCinco

\tituloI{Referencias Bibliográficas}

\begin{referencias}
\bibitem{tomic} TOMIĆ, N. Z. Economic model of microtransactions in video games. \textit{Journal of Economic Science Research}, 2019, 1 (1). DOI: 10.30564/jesr.v1i1.439
\bibitem{granma23} GRANMA. ¿Cómo los comercios estatales y privados pueden contratar los servicios de Código QR para pagos electrónicos? [En línea]. \textit{Granma}, 13 de agosto de 2023. [Consultado el: 23 de noviembre de 2026]. Disponible en: \url{https://www.granma.cu/cuba/2023-08-13/como-los-comercios-estatales-y-privados-pueden-contratar-los-servicios-de-codigo-qr-para-pagos-electronicos-13-08-2023-19-08-56}
\bibitem{etecsa} EMPRESA DE TELECOMUNICACIONES DE CUBA S.A. (ETECSA). Hitos de la historia. [En línea]. [Consultado el: 23 de noviembre de 2026]. Disponible en: \url{http://www.etecsa.cu/es/nosotros/hitos-de-la-historia}
\bibitem{cubadebate21} CUBADEBATE. Informatizar(nos): Puntos claves de la venta electrónica en el Comercio Interior. [En línea]. 5 de mayo de 2021. [Consultado el: 23 de noviembre de 2026]. Disponible en: \url{http://www.cubadebate.cu/noticias/2021/05/05/informatizar-nos-puntos-claves-de-la-venta-electronica-en-el-comercio-interior/}
\bibitem{mesa23b} MESA REDONDA. Pasarelas de pago e infraestructura tecnológica: Puntos claves para la bancarización. [En línea]. Cubadebate, 22 de agosto de 2023. [Consultado el: 23 de noviembre de 2026]. Disponible en: \url{http://mesaredonda.cubadebate.cu/mesa-redonda/2023/08/22/pasarelas-de-pago-e-infraestructura-tecnologica-puntos-claves-para-la-bancarizacion-video}
\bibitem{mesa23a} MESA REDONDA. Comercio electrónico en Cuba: Entre lo real y lo posible. [En línea]. Cubadebate, 6 de junio de 2023. [Consultado el: 23 de noviembre de 2026]. Disponible en: \url{http://mesaredonda.cubadebate.cu/mesa-redonda/2023/06/06/comercio-electronico-en-cuba-entre-lo-real-y-lo-posible-video}
\bibitem{oncuba23} REDACCIÓN OnCuba. Promueven en comercios estatales y privados contratación para pago electrónico. [En línea]. \textit{OnCubaNews}, 14 de agosto de 2023. [Consultado el: 23 de noviembre de 2026]. Disponible en: \url{https://oncubanews.com/cuba/promueven-en-comercios-estatales-y-privados-contratacion-para-pago-electronico}
\bibitem{rita24} RITA, P.; GUERREIRO, J.; RAMOS, R.; CAETANO, R. G. The role of microtransactions in impulse buying and purchase intention in the video game market. \textit{Entertainment Computing}, 2024, 50: 100693. DOI: 10.1016/j.entcom.2024.100693
\bibitem{hamarik17} HAMARI, J.; KERONEN, L. Why do people buy virtual goods: A meta-analysis. \textit{Computers in Human Behavior}, 2017, 71: p. 59-69. DOI: 10.1016/j.chb.2017.01.042
\bibitem{odlyzko03} ODLYZKO, A. The Case Against Micropayments. En: Wright, R. N. (editor). \textit{Financial Cryptography: 7th International Conference, FC 2003}. Berlín: Springer, 2003, p. 77-83. (Lecture Notes in Computer Science, vol. 2742). DOI: 10.1007/978-3-540-45126-6\_6
\bibitem{zendle20} ZENDLE, D.; MEYER, R.; BALLOU, N. The changing face of desktop video game monetisation: An exploration of exposure to loot boxes, pay to win, and cosmetic microtransactions in the most-played Steam games of 2010-2019. \textit{PLOS ONE}, 2020, 15 (5): e0232780. DOI: 10.1371/journal.pone.0232780
\bibitem{hamari15} HAMARI, J. Why do people buy virtual goods? Attitude toward virtual good purchases versus game enjoyment. \textit{International Journal of Information Management}, 2015, 35 (3): p. 299-308. DOI: 10.1016/j.ijinfomgt.2015.01.007
\bibitem{fielding00} FIELDING, R. T. \textit{Architectural Styles and the Design of Network-based Software Architectures}. Doctoral dissertation, University of California, Irvine, 2000. [Consultado el: 23 de noviembre de 2026]. Disponible en: \url{https://www.ics.uci.edu/~fielding/pubs/dissertation/top.htm}
\bibitem{stripelinks} STRIPE. Accept payments online without writing code (Payment Links). [En línea]. \textit{Stripe Documentation}. [Consultado el: 23 de noviembre de 2026]. Disponible en: \url{https://docs.stripe.com/payment-links}
\bibitem{rbi19} RESERVE BANK OF INDIA. \textit{Discussion Paper on Guidelines for Payment Gateways and Payment Aggregators}. Mumbai: Department of Payment and Settlement Systems, 2019. [Consultado el: 23 de noviembre de 2026]. Disponible en: \url{https://rbidocs.rbi.org.in/rdocs/PublicationReport/Pdfs/DPSSDISCUSSIONPAPEREFCF5B7E17F9431185BD4FD57E540F47.PDF}
\bibitem{stripeweb} STRIPE. Receive Stripe events in your webhook endpoint. [En línea]. \textit{Stripe Documentation}. [Consultado el: 23 de noviembre de 2026]. Disponible en: \url{https://docs.stripe.com/webhooks}
\bibitem{rfc2104} KRAWCZYK, H.; BELLARE, M.; CANETTI, R. \textit{HMAC: Keyed-Hashing for Message Authentication}. RFC 2104. Internet Engineering Task Force, 1997. [Consultado el: 23 de noviembre de 2026]. Disponible en: \url{https://datatracker.ietf.org/doc/html/rfc2104}
\bibitem{stripeidem} STRIPE. Idempotent requests. [En línea]. \textit{Stripe API Reference}. [Consultado el: 23 de noviembre de 2026]. Disponible en: \url{https://docs.stripe.com/api/idempotent_requests}
\bibitem{owasp23} OWASP FOUNDATION. \textit{OWASP API Security Top 10 - 2023}. [En línea]. 2023. [Consultado el: 23 de noviembre de 2026]. Disponible en: \url{https://owasp.org/API-Security/editions/2023/en/0x11-t10/}
\bibitem{etecsat} EMPRESA DE TELECOMUNICACIONES DE CUBA S.A. (ETECSA). Transfermóvil. [En línea]. [Consultado el: 23 de noviembre de 2026]. Disponible en: \url{http://www.etecsa.cu/es/personas/transfermovil}
\bibitem{xetid} XETID. ENZONA. [En línea]. \textit{Portal Web de la XETID}. [Consultado el: 23 de noviembre de 2026]. Disponible en: \url{https://www.xetid.cu/es/productos/enzona}
\bibitem{eltoque20} BOZA IBARRA, G. Pasarelas de pago en Cuba: ¿listas? [En línea]. \textit{elTOQUE}, 10 de febrero de 2020. [Consultado el: 23 de noviembre de 2026]. Disponible en: \url{https://eltoque.com/transfermovil-enzona-pasarelas-de-pago-en-cuba-ventajas-retos}
\bibitem{larman03} LARMAN, C.; BASILI, V. R. Iterative and incremental development: a brief history. \textit{Computer}, 2003, 36 (6): p. 47-56. DOI: 10.1109/MC.2003.1204375
\bibitem{typescript} MICROSOFT. \textit{TypeScript Documentation}. [En línea]. [Consultado el: 23 de noviembre de 2026]. Disponible en: \url{https://www.typescriptlang.org/docs/}
\bibitem{nextjs} VERCEL. \textit{Next.js Documentation}. [En línea]. [Consultado el: 23 de noviembre de 2026]. Disponible en: \url{https://nextjs.org/docs}
\bibitem{postgres} THE POSTGRESQL GLOBAL DEVELOPMENT GROUP. Transactions. En: \textit{PostgreSQL Documentation}. [En línea]. [Consultado el: 23 de noviembre de 2026]. Disponible en: \url{https://www.postgresql.org/docs/current/tutorial-transactions.html}
\bibitem{prisma} PRISMA. \textit{Prisma ORM Documentation}. [En línea]. [Consultado el: 23 de noviembre de 2026]. Disponible en: \url{https://www.prisma.io/docs/orm}
\bibitem{zod} ZOD. \textit{Zod: TypeScript-first schema validation}. [En línea]. [Consultado el: 23 de noviembre de 2026]. Disponible en: \url{https://zod.dev}
\bibitem{minio} MINIO, INC. \textit{MinIO Object Store Documentation}. [En línea]. [Consultado el: 23 de noviembre de 2026]. Disponible en: \url{https://docs.min.io/community/minio-object-store/index.html}
\bibitem{chacon14} CHACON, S.; STRAUB, B. \textit{Pro Git}. 2.ª ed. Apress, 2014. [En línea]. [Consultado el: 23 de noviembre de 2026]. Disponible en: \url{https://git-scm.com/book/en/v2}
\end{referencias}

% =====================================================================
%  Sección 11 — Anexos (Arial 11, interlineado sencillo)
% =====================================================================
\seccion{70.75bp}{2.55bp}{87.59bp}{28.77bp}{10.31bp}
\renewcommand\textoEncabezado{Anexos}
\doceUnoCinco

\tituloI[36bp]{Anexos}

\onceSencillo
<Contenido de los anexos con igual tipo de fuente Arial, pero a tamaño 11 puntos e interlineado 1.0 puntos. Debe tratar de sólo utilizarse aquellos anexos imprescindibles para complementar lo presentado en la memoria escrita y que no excedan las ocho (8) o diez (10 páginas). Deben aparecer uno a continuación del otro sin necesidad de saltos de página entre estos>

{\doceSencillo\vacio}

\end{document}
